% reality/topology_and_special_elements-DE.tex
\chapter{Topologie und besondere Elemente}
\label{chap:topology-and-special-elements}
\index{Topologie}
\index{Besondere Elemente!Eisenbahnnetz}

Eisenbahninfrastruktur ist vernetzt. Die Interpretation der Realität erfordert
daher neben dem geometrischen Zustand einen ausdrücklichen Topologiekontext.

\section{Topologiekontext}
\begin{definition}[Eisenbahnnetzkontext]
\index{Topologie!Eisenbahnnetzkontext}
Der Eisenbahnnetzkontext ist die Menge der Konnektivitätsbeziehungen zwischen
ausgerichteten Gleissegmenten, Verzweigungen, Zusammenführungen und
Routenfortsetzungen.
\end{definition}
Der Topologiekontext bestimmt Routenkontinuität und die Interpretation von
Stationierungs-, Projektions- und Beobachtungsdaten in mehrgleisigen Systemen.

\section{Besondere Elemente}
Weichen, Kreuzungen, Bahnhofsköpfe und Rangierstrukturen sind besondere
Elemente, an denen geometrische und operative Verzweigung ausdrücklich wird.
\begin{definition}[Verbindungspunkt]
\index{Topologie!Verbindungspunkt}
Ein Verbindungspunkt ist ein Ort, an dem zwei oder mehr Alignment-Segmente
topologisch verknüpft sind.
\end{definition}
Diese Elemente definieren konstruktive Alignment-Identität nicht neu; sie
stellen den Netzkontext bereit, in dem Identitäten interpretiert werden.

\section{Trennung von Geometrie und Topologie}
Eine nützliche Abstraktion trennt Geometrie von Konnektivität und ermöglicht
zugleich deren ausdrückliche Kopplung während der Interpretation.
\begin{definition}[Graphdarstellung]
\index{Topologie!Graphdarstellung}
Ein Eisenbahnnetz kann als Graph dargestellt werden, dessen Knoten
Verbindungspunkte und dessen Kanten Alignment-Segmente bezeichnen.
\end{definition}
Diese Darstellung verdeutlicht, dass Topologiekontext und geometrischer
Zustand verschieden, in Realitätsabläufen aber gemeinsam erforderlich sind.

\section{Datenmehrdeutigkeit und Inferenz}
Beobachtungsdatensätze spezifizieren die Topologie häufig unzureichend. Die
Inferenz von Routenkontinuität und Verzweigungszuordnung ist deshalb Teil der
Realitätsinterpretation und muss ausdrücklich behandelt werden.

\section{Verbindung zu AIM}
\begin{summary}
Topologie und besondere Elemente stellen den Netzkontext bereit, der für die
reale Interpretation gemessener und realisierter Eisenbahnzustände
erforderlich ist. In AIM wird Topologie als Kontextstruktur verwendet und
ersetzt nicht die konstruktive Alignment-Semantik.
\end{summary}
